\documentclass[11pt]{article}

\usepackage[margin=1in]{geometry}
\usepackage{graphicx}
\usepackage{amsmath}
\usepackage{booktabs}
\usepackage{longtable}
\usepackage{array}
\usepackage{hyperref}
\usepackage{xcolor}
\usepackage{enumitem}
\usepackage{microtype}
\usepackage[T1]{fontenc}
\usepackage[utf8]{inputenc}

\hypersetup{
  colorlinks=true,
  linkcolor=blue,
  citecolor=blue,
  urlcolor=blue
}

\title{DASHI: A Fail-Closed Constructive Physics-Unification Architecture\\
\large Carrier Geometry, Filtration, Projection-Defect Structure, and Mechanized Closure Frontiers}
\author{DASHI project}
\date{Draft of 2026-05-15}

\begin{document}
\sloppy
\maketitle

\noindent\fcolorbox{black!45}{black!3}{%
\begin{minipage}{\dimexpr\linewidth-2\fboxsep-2\fboxrule\relax}
\small
\textbf{Core formalism.}
Paper 1 uses the following grammar to stage DASHI's current
physics-unification carrier without treating the grammar itself as a universal
physical law:
\[
x \in \mathcal X \simeq (a,\nu,\tau,d,\ell,r),
\qquad
\tau \in \{-1,0,+1\}.
\]
Here \(a\) is a prime/supersingular-prime address, \(\nu\) is
\texttt{FactorVec}-style valuation data, \(\tau\) is the ternary coordinate,
\(d\) is refinement depth, \(\ell\) is the active lane, and \(r\) records
receipt/projection status. A surface \(s\) exposes only its admitted
projection:
\[
x = P_s(x) + D_s(x), \qquad
\rho(x,y)=\alpha^{k(x,y)},\; 0<\alpha<1.
\]
\(P_s(x)\) is the promoted component; \(D_s(x)\) is the residual defect; and
\(k(x,y)\) is the first incompatible address, valuation, or prefix depth.
Filtered operator surfaces use \(F_0 \subseteq F_1 \subseteq \cdots\),
\(\operatorname{gr}_nF=F_n/F_{n-1}\), and descended brackets only where
quotient, equivalence, norm, and bracket laws are inhabited. Lane actions
commute only on the official tracked independent \texttt{GL.FactorVec} route.
UFT/motif compression supplies reusable projected structure while residual
defects remain visible.
\end{minipage}}

\vspace{0.8em}

\begin{abstract}
DASHI is a constructive physics-unification program built around a typed
carrier geometry. Its central object is not a single fitted equation, but a
coherent mathematical substrate in which prime-coordinate carriers,
ultrametric refinement, projection-defect decomposition, filtered operator
surfaces, lane independence, local root-geometry boundaries, compression
semantics, and empirical-contact receipts can be staged inside one typed
architecture.

The manuscript uses a projection-defect schema to organize promoted content
and residual obstruction across those surfaces. Filtrations organize limiting
behavior, ultrametric addresses organize refinement depth, lane actions
organize independent coordinate sectors, and admissibility gates determine
which constructions may be read as physics-facing claims. Typed closure and
obstruction semantics are therefore not external project governance; they are
the enforcement layer of the unification substrate itself.

Paper 1 presents the current constructive frontier of this program. The
positive inventory includes direction-indexed G2 SFGC schema adoption, an
official tracked G6 cross-lane commuting route, constructive finite-support G3
selected-operator subtraction support, an E8/LILA integer and local semantic
boundary, a typed cross-domain projection-defect / MDL spine boundary, and an
inhabited bounded W3 below-Z Drell-Yan \texttt{t43} comparison-law receipt with
\texttt{chi2/dof = 2.1565191176} under threshold \texttt{4.0}. The contribution
is a mathematically coherent ultrametric carrier architecture for physics
unification whose closure semantics expose the exact frontier rather than hide
it.
\end{abstract}

\section{Introduction}

The starting point of DASHI is that a physics-unification program needs a
carrier rich enough to hold geometry, dynamics, compression, empirical contact,
and unresolved structure without forcing them into the same proof status. The
basic pattern is a state carried by prime/trit and \texttt{FactorVec}-style
coordinates, refined by ultrametric depth, acted on by lane-local operators,
and read through projections whose residuals remain visible. This is a
physics thesis before it is a governance thesis: the closure discipline is
needed because the proposed carrier is ambitious, not because the paper is
about process management.

The organizing picture is projection from structured possibility into a
surface-visible form. A candidate state can carry more information than a
particular physics-facing surface is allowed to expose. The promoted
projection is the piece that survives the current admissibility gate; the
defect records what is left unresolved, incompatible, or not yet transported
across the relevant boundary. Filtration gives limiting and contraction
language, ultrametric refinement gives address and precision language, and lane
independence gives compositionality between coordinate sectors.

This architecture lets several mathematical neighborhoods meet without being
collapsed into metaphor. Direction-indexed discrete sites supply local carrier
geometry. Finite-support operator surfaces supply controlled algebra for
filtered contraction targets. Tracked cross-lane commuting supplies a
coordinate-independence law for selected lane actions. Local E8/LILA surfaces
supply finite root-geometry and upstream-promotion boundaries. UFT-style
compression supplies refinement, motif, and residual language. The Drell-Yan
lane supplies a bounded example of empirical contact bound to a frozen typed
receipt.

\begin{figure}[ht]
  \centering
  \includegraphics[width=\linewidth]{Docs/PaperDraftWorkingFolder/figures/PublicPaperNarrativeFlow.pdf}
  \caption{Reader-facing narrative flow for Paper 1. The diagram foregrounds
  the constructive carrier substrate, local mathematical surfaces, empirical
  contact, compression geometry, and the later certification frontier. It is a
  narrative map, not a proof of completed unification.}
  \label{fig:narrative-flow}
\end{figure}

\section{Background and Related Work}

DASHI draws on standard mathematical neighborhoods including ultrametric and
p-adic structure, lattice and gauge language, contraction and associated
graded constructions, root-system classification, mechanized mathematics, and
compression-style trie geometry. The claims made here are restricted to the
repo-local receipt and obstruction surfaces cited in the claim ledger.

The formal-methods positioning is central. DASHI uses Agda as a typed
development surface: selected propositions, receipts, audits, and blockers are
represented as typed artifacts, and the absence of a required inhabitant is
manuscript-relevant state. This places the paper near proof assistants,
dependent type theory, mechanized mathematics, and proof-carrying artifacts,
while remaining narrower than an end-to-end certified physics pipeline.

Ultrametric and p-adic language organizes refinement by depth: nearby states
share high-level valuation or prefix structure, while separation can occur at
the first incompatible digit or coordinate. Hensel-style lifting is used only
as an external analogy for compatible refinement across precision levels. In
Paper 1 this is expository support for carrier and compression geometry, not
an independent proof of convergence.

\section{From Ternary Traversal to Carrier Geometry}

The origin path into DASHI is best read as an iterative formalization of
traversal: a state is first treated as a trit-valued position, then as a
refinable address, then as a carrier point whose projections and residuals are
tracked by typed operators. The point of the \texttt{3-6-9} or
\texttt{Base369} vocabulary is not numerological. Its safe role is to supply
cyclic finite carriers: triadic, hexadic, and nonary state spaces with
rotations, wrapped addition, and quotient/refinement seams. At the lowest
level, a trit distinguishes three statuses: positive, negative, and unresolved
or neutral.

Older dialectical language in the archive should be read in the same
disciplined way. A dialectical state may be carried through a transition while
retaining unresolved residue. A counter-state, conflict state, or overflow
state is not a mystical third force; it is a typed indication that the present
projection has not discharged all of the structure in the carrier. In the
current manuscript this becomes the projection-defect discipline: a selected
surface exposes \(P_s(x)\), while the residual \(D_s(x)\) records what is
blocked, incompatible, untransported, or not yet promoted.

The geometric intuition then moves from a one-dimensional cycle to traversal
through a cube-like state space. Rubik-like cube, hypercube, and voxel language
is useful only as intuition for finite-coordinate movement: each coordinate
records a local trit, phase, or valuation, and a traversal is a sequence of
admissible coordinate updates. Voxel language fits the same role at nonary
scale: a local \(3^2\) or \(3^k\) block can be read as a small cell of
structured state, and overflow means that the current cell cannot preserve the
relevant invariant without adding a refinement coordinate.

Video and state-continuity testing belong at this transition point as
engineering diagnostics. A video trace is a sampled path through state space:
continuity testing asks whether adjacent frames preserve the expected local
carrier relation, whether an apparent jump is an admissible carry/refinement,
and whether the projection defect increases at the boundary. Such tests can
motivate carrier choices and expose discontinuities, but they are not theorem
authority unless the trace, observation map, metric, and admissibility law are
bound by a repo-local receipt. The sibling \texttt{dashifine} and
\texttt{dashitest} repositories contain concrete video-side evidence at this
engineering level: GIF/frame visualizations for closure and ultrametric
diagnostics, codec and Vulkan/VAAPI video benchmark paths, live-sheet recording
scripts, and short MP4 smoke artifacts. These may be cited as sibling
diagnostic evidence, not as DASHI theorem authority, until ported into a local
receipt with command and artifact digests.

UFT/trie refinement gives this traversal a compression geometry. A trie node is
a coarse address; descending the trie adds precision. Nearby branches share
prefix structure, and separation occurs at the first incompatible depth. This
matches the ternary ultrametric reading already used in the carrier spine:
refinement depth is not just a storage index but a metric and semantic
relation. A motif is then a reusable local shape in the trie, not an equality
claim.

\texttt{FactorVec} is the corresponding prime-coordinate refinement of the
same idea. Instead of only trit positions, a state carries valuation data
across tracked prime lanes. Multi-lane transport, chamber legality, signed
scans, and defect summaries turn traversal into a controlled carrier geometry:
a move may be legal or illegal, stable or repatterning, contractive or
expansive, and these labels are computed at the representation layer rather
than inferred from prose.

This origin story gives a disciplined reading of the transition from trits to
carrier geometry. DASHI starts with ternary traversal because trits make
unresolvedness explicit. It passes through \texttt{Base369}, cube/voxel
traversal, video-like continuity traces, UFT tries, and \texttt{FactorVec}
lanes because each adds a more precise way to talk about transition, address,
refinement, and defect. None of this proves full physics unification,
automatic convergence, Monster representation semantics, empirical adequacy,
or domain recovery. It supplies a traceable formal path by which those stronger
claims would have to be typed, projected, tested, and closed.

\section{Why the Carrier Geometry Is Forced}

The carrier is not chosen by taste. It is the smallest structure that keeps
four pressures visible at once: unresolvedness, refinement, locality, and
projection. If a state can be affirmed, denied, or still undischarged by the
current surface, binary truth is already too small. A ternary coordinate is
the first stable carrier for that situation: \texttt{+1}, \texttt{-1}, and
\texttt{0} are not metaphysical labels, but bookkeeping for oriented evidence
and unresolved residue.

Once unresolved state is allowed to refine, a tree appears almost immediately.
A coarse node records what is visible at the current depth; descending the
trie adds distinctions. Two states that share a long prefix are close for any
observer or surface that cannot see past that depth, and the first
incompatible digit is the natural separation event. This is why the
ultrametric is not decoration. It is the metric forced by recursive
refinement: distance is controlled by the earliest depth at which two
histories, addresses, or motifs can no longer be jointly refined.

Prime-coordinate and \texttt{FactorVec} structure enter when refinement must
also remember independent axes of change. A factorized carrier records how
much of each tracked prime lane is present. If one update changes only the
\texttt{p} valuation and another changes only the \texttt{q} valuation, then
for \texttt{p != q} the two updates should commute on the tracked projection.
G6 is the formal version of that intuition: independent valuation coordinates
naturally commute, but only the official tracked theorem is allowed to say so
inside the paper.

Projection-defect decomposition is then inevitable rather than decorative.
Every observation, theorem surface, empirical receipt, or compression motif
exposes only part of the carrier. A less disciplined theory would forget the
remainder and later reintroduce it as an assumption. DASHI keeps the remainder
as \(D_s(x)\): the residual defect relative to a surface. This is why
obstruction belongs to the mathematics. It is the unpromoted part of the
carrier after projection, not a separate editorial apology.

Filtration is the same idea applied to operators. A local move is useful only
if it respects the depth and quotient structure on which the carrier is being
read. When an operation preserves the relevant filtration and descends to the
quotient, it can become a theorem-facing action on the associated graded
surface. When it does not, the failure is visible as a missing equivalence,
missing quotient carrier, or missing descended law. G3 is therefore not just a
blocker list; it is the exact place where finite-support operator intuition is
being asked to become quotient-correct algebra.

The root-geometry and empirical lanes fit the same pattern. E8/LILA supplies
finite local shape and symmetry boundaries where enumeration and promotion can
be audited. The Drell-Yan lane supplies a bounded observable projection where
a frozen comparison can be checked while stricter residual shape laws remain
visible. These are different projections of the same recursively refining
carrier grammar. Paper 1 is compelling only if the reader can see that unity
before seeing the blocker table.

\section{Core Constructive Architecture}

The core DASHI architecture is a carrier geometry before it is a list of
gates. The carrier is intended to hold prime-coordinate location, trit-valued
or unresolved state, \texttt{FactorVec}-style valuation data, lane-local
actions, depth/refinement information, and typed projections into the surfaces
a reader can inspect. This is the mathematical reason the architecture can
talk about discrete gauge sites, filtered operators, E8/LILA root boundaries,
compression motifs, and empirical comparison receipts without treating them as
unrelated subsystems.

\subsection{Carrier Coordinates and Depth}

The paper-facing carrier intuition is coordinate-rich but not continuously
overcommitted. Prime and supersingular-prime coordinates provide addressable
axes. Trit and Base369-style language supplies a compact way to discuss
positive, negative, and unresolved or neutral positions. \texttt{FactorVec}
surfaces make valuation information explicit enough for tracked lane actions.
UFT-style addresses add a refinement geometry: nearby addresses share
high-level prefix or valuation structure, and separation occurs at the first
incompatible depth.

This gives the carrier a natural physics reading. A state is not merely a
point in a set; it is a structured address with depth, projection grade,
valuation content, and possible unresolved branches. Refinement can therefore
be discussed without pretending that every refinement is already an observable
physical fact. Hensel-style language is useful as intuition for compatible
partial information lifting through increasing precision.

\subsection{Projection-Defect Decomposition}

The central local equation is the projection-defect split:
\[
  x = P(x) + D(x).
\]
Here \(P(x)\) denotes the promoted projection visible to a chosen surface, and
\(D(x)\) denotes the residual defect left by that projection. The point is not
that every domain has the same observables. The point is that DASHI repeatedly
uses the same structural move: expose the admissible part of a richer carrier,
preserve the residual, and make later promotion depend on a typed bridge
rather than a rhetorical analogy.

\subsection{Filtration, Lane Action, and Compression}

Filtration supplies the limiting direction of the architecture. It lets DASHI
ask whether a restricted operator surface descends to an associated graded
object, whether contraction language is quotient-correct, and whether norms or
brackets survive projection. Lane actions supply coordinate-local dynamics:
selected transformations can act on tracked prime sectors, and the official G6
theorem checks that selected cross-lane actions commute on the tracked route.

Compression is not decorative. The UFT and motif vocabulary supply a way to
talk about reusable structure, address refinement, branch compatibility, and
semantic distance. A motif is not equality; it is evidence of reusable
projection structure at a chosen grade. This is why compression belongs in the
physics paper: it is the language by which a rich carrier becomes searchable,
refinable, and composable without collapsing residual structure.

\subsection{Closure Semantics as Enforcement}

The closure layer enters after the carrier has something worth protecting.
Fail-closed promotion says that a local theorem, audit, or receipt promotes
only the claim authorized by its owning surface. A request surface can point to
an empirical, geometric, or upstream mathematical authority, but it does not
supply that authority. This is the enforcement architecture of the unification
program: it keeps the carrier ambitious while making the claim surface exact.

\section{Mathematical Constructions}

\begin{longtable}{p{0.12\linewidth}p{0.25\linewidth}p{0.28\linewidth}p{0.25\linewidth}}
\toprule
Label & Mathematical role & Current positive surface & Boundary \\
\midrule
G2 & Direction-indexed discrete carrier schema & SFGC schema adoption for current core-facing consumers & No Maxwell closure or field-equation recovery follows \\
G3 & Restricted finite-support operator filtration & Certified subtraction support for selected constructive finite-support operators & Associated-graded quotient infrastructure remains missing \\
G6 & Tracked cross-lane coordinate independence & Official tracked cross-lane commuting theorem and above-threshold independence & Old universal \texttt{LaneOperator} promotion remains blocked \\
E8/LILA & Local finite root-enumeration and upstream boundary & Integer/two-sparse completeness plus local semantic audit boundary & Upstream \texttt{E8RootEnumerationComplete} remains blocked \\
\bottomrule
\end{longtable}

\subsection{G2: Direction-Indexed Discrete Gauge Schema}

The G2 lane supplies the current direction-indexed carrier schema for local
discrete gauge geometry. The positive surface is
\texttt{canonicalG2SFGCGaugeFieldSchemaExtensionReceipt}, with the
no-remaining-missing surface
\texttt{noRemainingG2SFGCSchemaExtensionMissing}. The admissible paper claim is
narrow: direction-indexed SFGC schema adoption is complete for the current
core-facing consumers. This does not prove Maxwell closure, Yang-Mills
dynamics, field-equation recovery, empirical calibration, limited Standard
Model plus GR closure, GRQFT closure, or completed unification.

\subsection{G3: Selected Finite-Support Operators and the Quotient Gap}

The G3 lane is the restricted operator-calculus lane. The current positive
result is certified subtraction support for a selected constructive
finite-support subtype, through
\texttt{selectedFiniteSupportOperatorSubtractionCertifiedSupport},
\texttt{selectedFiniteSupportOperatorSubtractionSupportWitnessFromCertified},
and \texttt{G3SelectedSubtractionCertifiedSupport}. A genuine associated
graded layer still requires quotient classes of the form \texttt{F\_n /
F\_\{n-1\}} and the kernel-equivalence law:
\[
  [x] = [0] \quad \Longleftrightarrow \quad
  x \textrm{ lies in the denominator filtration piece.}
\]
Until the accepted equivalence, quotient carrier,
projection-respects-equivalence, and descended norm/bracket/isomorphism laws
are inhabited, Paper 1 cannot claim global Schrodinger promotion, arbitrary
operator finite support, norm-law closure, product support, bracket support,
or a completed Galilei contraction theorem.

\subsection{G6: Official Tracked Cross-Lane Commuting}

The G6 lane records cross-lane independence inside the prime-coordinate
carrier. An update to the \texttt{p} lane and an update to the \texttt{q} lane
factor through distinct valuation exponents, so for \texttt{p != q} the order
of those two tracked updates should not change the projected state. The
theorem matters because DASHI does not leave that independence as prose; it
pins it to the official tracked \texttt{GL.FactorVec} route. The positive
surfaces are
\texttt{canonicalG6OfficialTrackedCrossLaneCommutingTheorem} and
\texttt{canonicalG6OfficialTrackedAboveThresholdCoordinateIndependence}. The
old full \texttt{LaneOperator} interface is not promoted; its universal
vanished-prime identity law blocks the nontrivial scaling route used by the
tracked theorem. Paper 1 may claim the official tracked cross-lane commuting
theorem and tracked above-threshold coordinate independence, but not old
universal \texttt{LaneOperator} promotion, full G6 closure, or downstream
physics closure from G6 alone.

\subsection{E8/LILA: Local Completeness and Upstream Promotion Boundary}

The E8/LILA lane is a finite root-enumeration and promotion boundary, not a
completed E8-to-physics derivation. The positive and boundary surfaces are
\texttt{integerIndexedRootsCompleteForTwoSparseShapeTheorem},
\texttt{canonicalE8RootEnumerationCompletePromotionBoundary}, and
\texttt{canonicalE8UpstreamCompleteReceiptPromotionAudit}. Paper 1 may
describe local integer completeness and the exact upstream-promotion audit. It
must not claim upstream \texttt{E8RootEnumerationComplete}, E8-derived Standard
Model recovery, limited SM+GR, GRQFT closure, or completed unification from E8
structure.

\section{Empirical Contact}

The empirical-contact lane is intentionally narrow. Its role is to show that
the constructive architecture can bind a named physics observable to a frozen,
mechanically tracked receipt while still blocking stronger claims whose
authority, convention, or strict-fit obligations are missing.

The current positive surface is the bounded W3 below-Z Drell-Yan
comparison-law receipt over CMS-SMP-20-003 / HEPData \texttt{ins2079374/t43},
using the \texttt{ins2079374/t44} covariance convention. The receipt records
\texttt{chi2/dof = 2.1565191176}, threshold \texttt{4.0}, and mean
prediction/data \texttt{0.9941233097}. It has no above-Z or accepted-authority
promotion.

Two stricter protocols are explicitly negative. The adjacent-ratio receipt
records that the stricter \texttt{chi2/dof <= 2} falsification-lane target is
not satisfied. The strict log-covariance diagnostic is also negative:
\texttt{phiStarRatioPredictor} has strict log \texttt{chi2/dof =
283.45739523864586}, and \texttt{sigmaDashiV4Predictor} has strict log
\texttt{chi2/dof = 3180.211733150705}. The main obstruction is the log-linear
residual subspace; the next theorem target is a frozen no-refit route
discharging \texttt{StrictPassOrthogonalityObligation}, not a prose promotion.

The bounded W3 comparison-law receipt is not accepted W3 empirical authority.
No accepted provider token, non-postulated policy hook, or governance payload
currently inhabits \texttt{W3AcceptedEvidenceAuthorityToken}. W4 and W5 remain
blocked by missing accepted Drell-Yan luminosity/convention authority.

\section{Compression, Semantic Geometry, and Cross-Domain Spine}

The UFT vocabulary supplies a compression-facing view of the carrier:
addresses refine ultrametrically, motif membership is graded by projection
evidence, and unresolved or incompatible branches remain represented rather
than collapsed. The visual intuition is a floating trie of possible
refinements. A coarse address names a region of structure; descending the trie
adds precision, distinguishes branches, and exposes incompatibilities. Unlike
an ordinary lookup tree, the UFT reading is semantic and typed: branch distance
is ultrametric, branch membership is projection-grade dependent, and
unresolved branches can remain live without being forced into a binary
collapse.

Motifs then become reusable local shapes inside that trie. A motif does not
assert identity between two states. It says that two carrier regions share
enough projection structure to support a typed comparison at a chosen grade.
Exact members, compatible members, blocked members, and residual members are
different mathematical statuses. This lets the architecture reuse structure
without treating resemblance as theoremhood.

UFT-C, lattice-logic, null/hinge codes, SWAR implementation machinery,
Base369/trit notation, Hensel-style refinement, and Markov-after-quotient
readings are routed as background or appendix material unless a repo-local
surface promotes a stronger claim.

The cross-domain material is included as a theorem-target boundary, not as a
new promoted result. The current Agda surface
\texttt{canonicalCrossDomainVariationalSpineBoundary} records a common typed
object \texttt{(X, delta, pi, defect, gate, observation, symmetry)} shared by
physics, molecular, bonding, biological, and perceptual target rows. Its value
for Paper 1 is structural: it shows how projection, defect, admissible
descent, observation quotient, and symmetry-orbit language can be stated in
one typed grammar. It does not imply chemistry closure, bonding-energy
prediction, molecular spectra, protein folding, biological prediction,
perceptual empirical fit, or universal TOE closure.

\section{Obstruction Frontier}

The preceding sections put the carrier first: coordinate structure,
projection-defect decomposition, filtration, lane action, compression, local
mathematical surfaces, empirical contact, and cross-domain spine. The
obstruction frontier now records the exact residuals of that architecture. It
is not the emotional center of the paper, but it is the certification layer
that makes the physics-unification thesis admissible.

\begin{longtable}{p{0.16\linewidth}p{0.29\linewidth}p{0.27\linewidth}p{0.18\linewidth}}
\toprule
Gate & Current surface & First missing primitive/token & Forbidden promotion \\
\midrule
G2 & Direction-indexed SFGC schema adoption & Maxwell/curvature/field-equation carrier and receipts & Maxwell closure, physical curvature \\
G3 & Selected finite-support subtraction support & Quotient carrier, kernel-equivalence law, descended laws & Global operator or Schrodinger closure \\
G6 & Official tracked commuting route & Migration or split of old universal law & Old full \texttt{LaneOperator} promotion \\
E8/LILA & Local integer/root boundary & Upstream completion constructor/API & Upstream E8 completion \\
W3 & Bounded Drell-Yan receipt & \texttt{W3AcceptedEvidenceAuthorityToken} & Accepted empirical authority \\
W4/W5 & Request/diagnostic surfaces & Accepted DY convention and correction receipts & Matter field, stress-energy, sourced physics \\
GR & Future adapter path & Non-flat connection, curvature, Ricci, sourced interface & Einstein tensor, Schwarzschild, continuum GR \\
Cross-domain spine & Shared projection-defect / MDL schema & Calibration, universality, domain receipts & Chemistry/biology/perception recovery \\
Provenance & JMD/SSP/archive/source routing & Repo-local receipts or accepted authority & Provenance-as-proof \\
\bottomrule
\end{longtable}

The frontier has a constructive interpretation. Every row identifies a
residual defect: the exact bridge, authority, quotient law, adapter, or
calibration object that would be needed for a stronger projection. This is how
DASHI avoids converting an elegant carrier into an overclaimed physics result.

\section{Reproducibility}

The canonical authoring surface is
\texttt{Docs/PaperDraftWorkingFolder/Paper1\_Manuscript.md}. The claim-control
surface is \texttt{Docs/PaperDraftWorkingFolder/ClaimLedger.md}, and
source/provenance routing is recorded in
\texttt{Docs/PaperDraftWorkingFolder/SourceLedger.md},
\texttt{ArchiveTopicIntegrationPlan.md}, and
\texttt{ArchiveRelativeLocationReference.md}.

The narrow verification surface for this draft is:
\begin{verbatim}
plantuml -checkonly Docs/*.puml
./scripts/render_docs_diagrams.sh
git diff --check -- Docs/PaperDraftWorkingFolder Docs README.md
\end{verbatim}

\section{What This Paper Does Not Claim}

This paper does not claim that DASHI has completed physics unification. It does
not claim an accepted theory of everything, accepted external empirical
authority, physical calibration, sourced Einstein equations, non-flat GR
closure, Schwarzschild recovery, continuum recovery, GRQFT closure, limited
SM+GR closure, strict Drell-Yan success, above-Z Drell-Yan promotion, W4/W5
authority closure, upstream E8 completion, universal G6 closure, arbitrary
wavefunction-operator closure, chemistry closure, biology prediction,
perceptual validation, or cross-domain recovery equivalence.

\section{Conclusion}

DASHI should be read as a constructive physics-unification program centered on
a coherent carrier geometry. Prime-coordinate structure, ultrametric
refinement, projection-defect decomposition, filtered operator surfaces,
tracked lane actions, local root-geometry boundaries, semantic compression,
and bounded empirical receipts are staged as parts of one mathematical
architecture. The paper's thesis is that this architecture makes unification a
constructive closure problem rather than a sequence of informal analogies.

The current frontier is substantial but not final. G2 supplies a
direction-indexed carrier schema for current core-facing consumers. G3
supplies selected finite-support subtraction support and exposes the
associated-graded quotient gap. G6 supplies the official tracked cross-lane
commuting route. E8/LILA supplies local integer/root-boundary structure while
preserving upstream ownership. UFT and motif language supply the compression
geometry that keeps refinement, semantic distance, and residual structure in
the same typed vocabulary. The Drell-Yan lane supplies bounded empirical
contact and the stricter residual shape law that still fails.

Typed closure is what keeps this architecture from becoming another
unification narrative closed by assertion. It does not reduce the physics
ambition; it makes the ambition legible. The positive surfaces stand where
they are inhabited, and the stronger empirical, GR, upstream E8,
universal-operator, and completed-unification readings remain exact residual
obligations. Paper 1's contribution is therefore a mathematically coherent
ultrametric carrier architecture for physics unification with a visible,
mechanized closure frontier.

\nocite{dashi-claim-ledger-2026,dashi-source-ledger-2026,dashi-public-narrative-flow-2026,agda-project,hepdata-placeholder}
\bibliographystyle{plain}
\bibliography{Docs/PaperDraftWorkingFolder/Paper1_References}

\end{document}
